\section{第16章 コンピューティング基礎}
原典：Guide to the Software Engineering Body of Knowledge (SWEBOK Guide) v4.0a, Chapter 16 Computing Foundations。形式参考：01章 ソフトウェア要求の日本語まとめ。

\begin{pointbox}{この資料の主張}
コンピューティング基礎は、ソフトウェア技術者が設計・実装・性能・信頼性・安全性の判断をするための土台である。プログラムを書く能力だけでは、どのデータ構造を選ぶか、どの計算量が許容できるか、どの基本ソフトウェアやデータベースやネットワーク構成を選ぶかを説明できない。本章は、コンピュータの構成、データ構造とアルゴリズム、プログラミング言語、基本ソフトウェア、データベース、ネットワーク、人間要因、人工知能・機械学習を、ソフトウェア工学の判断材料として整理する章である。

\end{pointbox}

\begin{pointbox}{この資料の読み方}
専門用語は原則として日本語で示し、元の英語名を確認したい語は括弧内に併記する。初学者が前提知識なしに読めるように、各節では定義、実務での意味、判断の観点を分けて説明する。本資料は公式訳ではなく、SWEBOK 16章の内容を学習用に再構成した要約である。

\end{pointbox}
